% Options for packages loaded elsewhere
\PassOptionsToPackage{unicode}{hyperref}
\PassOptionsToPackage{hyphens}{url}
\PassOptionsToPackage{space}{xeCJK}
\documentclass[
  a4paper,11pt]{ltjsarticle}
\usepackage{xcolor}
\usepackage[margin=1in]{geometry}
\usepackage{amsmath,amssymb}
\usepackage{graphicx}
\setcounter{secnumdepth}{-\maxdimen} % remove section numbering
\usepackage{iftex}
\ifPDFTeX
  \usepackage[T1]{fontenc}
  \usepackage[utf8]{inputenc}
  \usepackage{textcomp}
\else
  \usepackage{unicode-math}
  \defaultfontfeatures{Scale=MatchLowercase}
  \defaultfontfeatures[\rmfamily]{Ligatures=TeX,Scale=1}
\fi
\usepackage{lmodern}
\ifPDFTeX\else
  \setmainfont[]{Times New Roman}
  \ifXeTeX
    \usepackage{xeCJK}
    \setCJKmainfont[]{Hiragino Mincho ProN}
  \fi
  \ifLuaTeX
    \usepackage[]{luatexja-fontspec}
    \setmainjfont[]{Hiragino Mincho ProN}
  \fi
\fi
\IfFileExists{upquote.sty}{\usepackage{upquote}}{}
\IfFileExists{microtype.sty}{%
  \usepackage[]{microtype}
  \UseMicrotypeSet[protrusion]{basicmath}
}{}
\makeatletter
\@ifundefined{KOMAClassName}{%
  \IfFileExists{parskip.sty}{%
    \usepackage{parskip}
  }{%
    \setlength{\parindent}{0pt}
    \setlength{\parskip}{6pt plus 2pt minus 1pt}}
}{% if KOMA class
  \KOMAoptions{parskip=half}}
\makeatother
\usepackage{longtable,booktabs,array}
\newcounter{none}
\usepackage{calc}
\usepackage{etoolbox}
\makeatletter
\patchcmd\longtable{\par}{\if@noskipsec\mbox{}\fi\par}{}{}
\makeatother
\IfFileExists{footnotehyper.sty}{\usepackage{footnotehyper}}{\usepackage{footnote}}
\makesavenoteenv{longtable}
\usepackage{bookmark}
\IfFileExists{xurl.sty}{\usepackage{xurl}}{}
\urlstyle{same}
\usepackage{hyperref}
\hypersetup{
  hidelinks,
  pdfcreator={LaTeX via pandoc}}
\setlength{\emergencystretch}{3em}
\providecommand{\tightlist}{%
  \setlength{\itemsep}{0pt}\setlength{\parskip}{0pt}}
\usepackage{amsthm}
\newtheorem{theorem}{定理}[section]
\newtheorem{proposition}[theorem]{命題}
\newtheorem{lemma}[theorem]{補題}
\newtheorem{corollary}[theorem]{系}
\theoremstyle{definition}
\newtheorem{definition}[theorem]{定義}
\theoremstyle{remark}
\newtheorem{remark}[theorem]{注意}
\newtheorem{observation}[theorem]{観察}
\begin{document}
\section{6次元超直方体の集合構造とその組合せ論的性質}\label{ux6b21ux5143ux8d85ux76f4ux65b9ux4f53ux306eux96c6ux5408ux69cbux9020ux3068ux305dux306eux7d44ux5408ux305bux8ad6ux7684ux6027ux8cea}

\subsection{------ スピン分類・集合の数え上げ・符号積
------}\label{ux30b9ux30d4ux30f3ux5206ux985eux96c6ux5408ux306eux6570ux3048ux4e0aux3052ux7b26ux53f7ux7a4d}

\textbf{著者}: 木原 範昭 (WF System Co., Ltd.) \textbf{日付}: 2026年4月
\textbf{DOI}: 10.5281/zenodo.19646652

\begin{center}\rule{0.5\linewidth}{0.5pt}\end{center}

\subsection{0. 本稿の目的}\label{ux672cux7a3fux306eux76eeux7684}

本稿は、5 次元超直方体に離散ラベル軸 \(Q\) を追加した 6
次元構造を定義し、その集合構造から導かれる組合せ論的性質を述べる。

本体（§1--§8）では、本モデルの数学的構造と、そこから純粋に組合せ論的に導出される帰結のみを主張する。標準模型との対応は
§9
で「解釈の一例」として提示するが、これは本論文の主張ではない。モデルの集合構造が標準模型の粒子分類とどのように対応付けうるかを示す一例に過ぎない。

\textbf{本モデルの外部仮定について}: 本稿は 5
次元超直方体を出発点とする。なぜ 5 次元であるか、なぜ 5
軸のうち特定の軸を後に「空間軸」として扱うかは、\textbf{本モデルの外部仮定}であり、モデル内部では正当化されない。観測される時空が
3 次元空間 + 1
次元時間であるという経験的事実への対応として、この次元設定を採用している。

\begin{center}\rule{0.5\linewidth}{0.5pt}\end{center}

\subsection{1. 6
次元構造の定義}\label{ux6b21ux5143ux69cbux9020ux306eux5b9aux7fa9}

\subsubsection{1.1 定義}\label{ux5b9aux7fa9}

本稿で扱う 6 次元構造は以下で定義される。

\begin{itemize}
\tightlist
\item
  \textbf{幾何学的構造}: 5
  次元超直方体の組合せ構造（頂点、辺、面、3-cell、4-facet、5-cell 本体）
\item
  \textbf{ラベル軸}: 4 値の離散ラベル軸 \(Q \in \{0, 1, 2, 3\}\)
\end{itemize}

5 次元超直方体の頂点座標は \((x_1, x_2, x_3, x_4, x_5)\)
で表記し、各座標は\textbf{任意の実数値をとる}。特定の命題が集合構造において座標値に条件を課す場合は、その命題の中で局所的に宣言する。

第 6 軸 \(Q\)
は幾何学的な座標ではなく、各状態に付随する離散ラベル値である。値域は
\(\{0, 1, 2, 3\}\) の 4 値とする。この 4 値は 2
ビットの構造を持ち、\(Q = 2c_2 + c_3\)（\(c_2, c_3 \in \{0, 1\}\)）として符号化される。\(Q = 0\)
は Q 軸成分が零であること、すなわち Q
軸に対して直交する状態を意味する。\(Q\)
軸の追加によって超直方体の組合せ構造（頂点数、辺数、面数など）は変化しない。

\subsubsection{\texorpdfstring{1.2 \(k\)
次元超直方体の要素数}{1.2 k 次元超直方体の要素数}}\label{k-ux6b21ux5143ux8d85ux76f4ux65b9ux4f53ux306eux8981ux7d20ux6570}

\(k\) 次元超直方体の \(j\) 次元面（\(j\)-face）の個数を \(f_j(k)\)
とする。各要素数を低次元から書き下すことで、一般式を読者が検算できる形で提示する。

\textbf{頂点数 \(f_0(k)\)}: 各軸について両端の 2
状態を区別する組合せ数として

{\def\LTcaptype{none} % do not increment counter
\begin{longtable}[]{@{}ccccccc@{}}
\toprule\noalign{}
\(k\) & 0 & 1 & 2 & 3 & 4 & 5 \\
\midrule\noalign{}
\endhead
\bottomrule\noalign{}
\endlastfoot
\(f_0(k)\) & 1 & 2 & 4 & 8 & 16 & 32 \\
\end{longtable}
}

\[f_0(k) = 2^k\]

\textbf{辺の数 \(f_1(k)\)}: 1 本の軸方向に伸び、残り \((k-1)\)
本の軸の両端位置で決まる

{\def\LTcaptype{none} % do not increment counter
\begin{longtable}[]{@{}ccccccc@{}}
\toprule\noalign{}
\(k\) & 0 & 1 & 2 & 3 & 4 & 5 \\
\midrule\noalign{}
\endhead
\bottomrule\noalign{}
\endlastfoot
\(f_1(k)\) & 0 & 1 & 4 & 12 & 32 & 80 \\
\end{longtable}
}

\[f_1(k) = k \cdot 2^{k-1}\]

検算: \(k=3\) で \(3 \cdot 4 = 12\)（立方体の辺 12 本）、\(k=5\) で
\(5 \cdot 16 = 80\)。

\textbf{正方形面の数 \(f_2(k)\)}: 2 本の軸で張られる面

{\def\LTcaptype{none} % do not increment counter
\begin{longtable}[]{@{}ccccccc@{}}
\toprule\noalign{}
\(k\) & 0 & 1 & 2 & 3 & 4 & 5 \\
\midrule\noalign{}
\endhead
\bottomrule\noalign{}
\endlastfoot
\(f_2(k)\) & 0 & 0 & 1 & 6 & 24 & 80 \\
\end{longtable}
}

\[f_2(k) = \binom{k}{2} \cdot 2^{k-2}\]

検算: \(k=3\) で \(3 \cdot 2 = 6\)（立方体の面 6 枚）、\(k=5\) で
\(10 \cdot 8 = 80\)。

\textbf{立方体胞（3-cell）の数 \(f_3(k)\)}: 3 本の軸で張られる立方体胞

{\def\LTcaptype{none} % do not increment counter
\begin{longtable}[]{@{}ccccccc@{}}
\toprule\noalign{}
\(k\) & 0 & 1 & 2 & 3 & 4 & 5 \\
\midrule\noalign{}
\endhead
\bottomrule\noalign{}
\endlastfoot
\(f_3(k)\) & 0 & 0 & 0 & 1 & 8 & 40 \\
\end{longtable}
}

\[f_3(k) = \binom{k}{3} \cdot 2^{k-3}\]

検算: \(k=4\) で \(4 \cdot 2 = 8\)、\(k=5\) で \(10 \cdot 4 = 40\)。

\textbf{4-facet の数 \(f_4(k)\)}: 4 本の軸で張られる 4 次元セル

{\def\LTcaptype{none} % do not increment counter
\begin{longtable}[]{@{}ccccccc@{}}
\toprule\noalign{}
\(k\) & 0 & 1 & 2 & 3 & 4 & 5 \\
\midrule\noalign{}
\endhead
\bottomrule\noalign{}
\endlastfoot
\(f_4(k)\) & 0 & 0 & 0 & 0 & 1 & 10 \\
\end{longtable}
}

\[f_4(k) = \binom{k}{4} \cdot 2^{k-4}\]

検算: \(k=5\) で \(5 \cdot 2 = 10\)。

\textbf{5-cell 本体の数 \(f_5(k)\)}: \(k=5\) の場合

\[f_5(5) = \binom{5}{5} \cdot 2^0 = 1\]

\subsubsection{一般式}\label{ux4e00ux822cux5f0f}

以上をまとめると、\(k\) 次元超直方体の \(j\) 次元面の個数は

\[f_j(k) = \binom{k}{j} \cdot 2^{k-j} \qquad (0 \leq j \leq k)\]

で与えられる。これは、\(k\) 本の軸から \(j\) 本を選んで \(j\)-face
を張り（\(\binom{k}{j}\) 通り）、残り \((k-j)\)
本の軸について各両端のどちらに位置するかを指定する（\(2^{k-j}\)
通り）ことに対応する。

\textbf{検算（オイラー標数）}:
\(\sum_{j=0}^{k} (-1)^j f_j(k) = (2-1)^k = 1\)
により、超直方体が可縮空間であることと整合する。

\subsubsection{\texorpdfstring{本稿で用いる値（\(k=5\)）}{本稿で用いる値（k=5）}}\label{ux672cux7a3fux3067ux7528ux3044ux308bux5024k5}

{\def\LTcaptype{none} % do not increment counter
\begin{longtable}[]{@{}clc@{}}
\toprule\noalign{}
\(j\) & 要素 & \(f_j(5)\) \\
\midrule\noalign{}
\endhead
\bottomrule\noalign{}
\endlastfoot
0 & 頂点 & 32 \\
1 & 辺 & 80 \\
2 & 正方形面 & 80 \\
3 & 立方体胞（3-cell） & 40 \\
4 & 4-facet & 10 \\
5 & 5-cell 本体 & 1 \\
\end{longtable}
}

\subsubsection{1.3
軸の対称性と集合符号の定義}\label{ux8ef8ux306eux5bfeux79f0ux6027ux3068ux96c6ux5408ux7b26ux53f7ux306eux5b9aux7fa9}

\paragraph{1.3.1 5
軸の完全対称性}\label{ux8ef8ux306eux5b8cux5168ux5bfeux79f0ux6027}

\textbf{命題 1.1（5 軸の完全対称性）}. 5 次元超直方体の 5
本の幾何学的軸は、組合せ構造上すべて対等である。任意の軸の置換は超直方体の自己同型写像を与え、§1.2
の要素数 \(f_j(k)\) および後続の幾何学的命題に関して 5
軸は区別されない。

\textbf{証明}. \(k\) 次元超直方体の組合せ構造の自己同型群は超八面体群
\(B_k\) であり、任意の軸置換を含む。したがって 5
軸は組合せ的に完全対称である。\(\square\)

\paragraph{1.3.2 軸の分類}\label{ux8ef8ux306eux5206ux985e}

\textbf{定義 1.2（軸の表記ラベル）}. 5 本の軸
\(x_1, x_2, x_3, x_4, x_5\) に対して以下の表記ラベルを与える。

\begin{itemize}
\tightlist
\item
  \(x_1, x_2, x_3\) を \(x, y, z\) と表記する
\item
  \(x_4\) を \(t\) と表記する
\item
  \(x_5\) を \(R\) と表記する
\end{itemize}

命題 1.1 により 5
軸は組合せ構造上対等であるから、どの軸にどのラベルを付すかは定義の選択に過ぎない。これらのラベルは集合符号の定義（§1.3.3）およびスピンの定義（§2）の表記の前提として用いる。

\textbf{\(x, y, z, t, R\)
はあくまで表記上のラベルであり、幾何学的軸としての区別を意味しない。}
物理的解釈を各軸に割り当てる操作は
§9（解釈の一例）で行うが、これはモデル外部からの解釈であり、§1
の幾何学的構造とは独立である。

\paragraph{1.3.3
集合符号の厳密な定義}\label{ux96c6ux5408ux7b26ux53f7ux306eux53b3ux5bc6ux306aux5b9aux7fa9}

各幾何学的軸の成分
\(x_i\)（\(i = 1, \ldots, 5\)）は、その軸方向の座標値が満たすべき\textbf{範囲制約}を表す記号である。本稿では以下の
5 種類の記号を用いる。

{\def\LTcaptype{none} % do not increment counter
\begin{longtable}[]{@{}cll@{}}
\toprule\noalign{}
記号 & 意味 & 座標値の範囲 \\
\midrule\noalign{}
\endhead
\bottomrule\noalign{}
\endlastfoot
\(+\) & 正値 & \(x_i > 0\) \\
\(-\) & 負値 & \(x_i < 0\) \\
\(0\) & 原点 & \(x_i = 0\) \\
\(\pm\) & 非零値 & \(x_i \neq 0\)（正または負） \\
\(*\) & 任意 & \(x_i\) は任意の実数 \\
\end{longtable}
}

\(Q\) 軸についても同様に、値域を表す記号を用いる。

{\def\LTcaptype{none} % do not increment counter
\begin{longtable}[]{@{}
  >{\centering\arraybackslash}p{(\linewidth - 4\tabcolsep) * \real{0.1343}}
  >{\raggedright\arraybackslash}p{(\linewidth - 4\tabcolsep) * \real{0.2090}}
  >{\raggedright\arraybackslash}p{(\linewidth - 4\tabcolsep) * \real{0.6567}}@{}}
\toprule\noalign{}
\begin{minipage}[b]{\linewidth}\centering
記号
\end{minipage} & \begin{minipage}[b]{\linewidth}\raggedright
意味
\end{minipage} & \begin{minipage}[b]{\linewidth}\raggedright
\(Q\) の値
\end{minipage} \\
\midrule\noalign{}
\endhead
\bottomrule\noalign{}
\endlastfoot
\(0\) & 零ラベル（Q 軸に直交） & \(Q = 0\) \\
\(1\)--\(3\) & 指定離散値 & \(Q\) は指定された整数値 \\
\(*\) & 任意 & \(Q \in \{0, 1, 2, 3\}\) のいずれか \\
\end{longtable}
}

\(Q\) の 4 値は 2 ビット構造 \((c_2, c_3)\)
を持つ。\(Q = 0\)（\(c_2 = c_3 = 0\)）は \(Q\) 軸成分が零であり、\(Q\)
軸に対して直交する状態を表す。\(Q \neq 0\) の 3
状態（\(c_2 c_3 \in \{01, 10, 11\}\)）は \(Q\)
軸上の非零成分を持つ状態である。物理的解釈は §9.2 で扱う。

\textbf{注記}.
集合符号は座標値そのものではなく、座標値が満たすべき\textbf{条件}を指定する記号である。§1.1
で述べた通り、本モデルの座標値は任意の実数値をとり得るが、特定の命題や集合構造においては、該当する軸の座標値に上記の範囲制約を課す。

\textbf{例}. 集合 \((+, +, 0, 0, 0, 0)\)
は、\(x_1 > 0, x_2 > 0, x_3 = 0, x_4 = 0, x_5 = 0, Q = 0\)
を満たす状態の集合を指す。

\begin{center}\rule{0.5\linewidth}{0.5pt}\end{center}

\subsection{2. スピンの定義}\label{ux30b9ux30d4ux30f3ux306eux5b9aux7fa9}

\textbf{定義 2.1（スピン）}. 集合 \((x, y, z, t, R, Q)\)
に対して、スピン \(s\) を以下で定義する。

\textbf{フェルミオン/ボソンの判定}: \(xyzt\) の 4 軸のうち、集合符号が
\(0\) でない軸の本数を \(k\) とする。

\begin{itemize}
\tightlist
\item
  \(k = 2\) のとき: \textbf{フェルミオン型}
\item
  \(k \neq 2\) のとき: \textbf{ボソン型}
\end{itemize}

\textbf{スピン量子数の決定}: スピンに寄与する軸の本数 \(n\)
を以下で定める。\(Q\) 軸については、\(Q\)
ラベル間遷移（§3）のみがスピンに寄与し、静的な \(Q\)
ラベル値（\(Q \neq 0\) であっても遷移を伴わない場合）は \(n\)
に算入しない。

\begin{itemize}
\tightlist
\item
  フェルミオン型: \(n\) は \(\{R,\, Q\text{-遷移}\}\)
  のうち非零の本数。\(s = n + 1/2\)
\item
  ボソン型: \(n\) は \(\{t,\, R,\, Q\text{-遷移}\}\)
  のうち非零の本数。\(s = n\)
\end{itemize}

\(n\) のとりうる値はフェルミオン型で \(0, 1, 2\)、ボソン型で
\(0, 1, 2, 3\) であるから、本モデルにおけるスピン \(s\)
のとりうる値は以下となる。

{\def\LTcaptype{none} % do not increment counter
\begin{longtable}[]{@{}ccc@{}}
\toprule\noalign{}
\(n\) & ボソン \(s = n\) & フェルミオン \(s = n + 1/2\) \\
\midrule\noalign{}
\endhead
\bottomrule\noalign{}
\endlastfoot
0 & \(s = 0\) & \(s = 1/2\) \\
1 & \(s = 1\) & \(s = 3/2\) \\
2 & \(s = 2\) & \(s = 5/2\) \\
3 & \(s = 3\) & ------ \\
\end{longtable}
}

\textbf{注記 1}. \(xyzt\) の 4
軸は初期対称性において完全に対等であり、そのうち 1
軸が「時間的」として分離する過程は自発的対称性の破れとして解釈される（§9.1）。フェルミオン判定条件
\(k = 2\) は、4 つの対等な軸から 2 つが非零となる場合に対応し、\(xyz\)
の 3 軸に限定しない。

\textbf{注記 2}.
本定義は本モデル内部の組合せ論的な定義である。ここで「ボソン」「フェルミオン」と呼ぶのは本モデル内の用語であり、物理的なスピン角運動量、統計性、\((-1)^{2s}\)
変換性などとの対応は §9（解釈の一例）で扱う。

\textbf{注記 3}. フェルミオン型とボソン型で \(n\)
に算入する軸が異なる（フェルミオンは \(t\)
を算入しない）。この非対称性は以下の幾何学的動機に基づく。フェルミオン型では
\(xyzt\) の 4 軸のうちちょうど 2 本が非零であり、この 2 本の選択がすでに
\(\binom{4}{2} = 6\) 通りの自由度を持つ。このうち \(t\)
軸を含む組合せ（3 通り）と含まない組合せ（3 通り）は、\(k = 2\)
の判定で共に参加済みである。したがって \(t\)
の自由度はフェルミオン性の決定に消費されており、スピン量子数
\(s = n + 1/2\) の \(1/2\)
に反映される。数学的に等価な表記として、\(\{t, R, Q\text{-遷移}\}\)
の非零本数 \(n'\) を用いて \(s = (n' - 1) + 1/2 = n' - 1/2\)
とも書ける（\(t\) が非零のフェルミオンの場合
\(n' = n + 1\)）が、本稿では定義の見通しの良さから上記の非対称形を採用する。

\begin{center}\rule{0.5\linewidth}{0.5pt}\end{center}

\subsection{\texorpdfstring{3. \(Q\)
ラベル間遷移の数え上げ}{3. Q ラベル間遷移の数え上げ}}\label{q-ux30e9ux30d9ux30ebux9593ux9077ux79fbux306eux6570ux3048ux4e0aux3052}

\(Q\) 軸の 4 値のうち非零の 3 値 \(\{1, 2, 3\}\)
の間の遷移を考える。\(Q = 0\) は \(Q\) 軸に直交する状態であり、\(Q\)
軸上の遷移には関与しない。遷移は \((Q_\text{初期}, Q_\text{終期})\)
の順序対で記述される。

2 ビット構造 \((c_2, c_3)\) に基づき、非零値
\(c_2 c_3 \in \{01, 10, 11\}\) の 3 値間の順序対は \(3 \times 3 = 9\)
通りである（対角成分 3 通りを含む）。

\textbf{注記}. 9
通りから独立な遷移状態を抽出する操作（例えば対角成分の線形結合の除去）は、表現論的な構造に依存するため、本体では行わない。独立状態数の物理的解釈は
§9.3 で扱う。

\begin{center}\rule{0.5\linewidth}{0.5pt}\end{center}

\subsection{4. スピン 1
集合の数え上げ}\label{ux30b9ux30d4ux30f3-1-ux96c6ux5408ux306eux6570ux3048ux4e0aux3052}

スピン 1 の集合は定義 2.1 により、ボソン型（\(xyzt\) のうち非零軸が
\(k \neq 2\)）かつ \(n = 1\) の状態である。\(xyzt\) 軸がすべて
\(0\)（\(k = 0\)）の場合を考える。

\(\{t, R, Q\text{-遷移}\}\) のうち非零が 1
本の状態は以下のように分類される。

{\def\LTcaptype{none} % do not increment counter
\begin{longtable}[]{@{}
  >{\raggedright\arraybackslash}p{(\linewidth - 4\tabcolsep) * \real{0.2105}}
  >{\raggedright\arraybackslash}p{(\linewidth - 4\tabcolsep) * \real{0.5263}}
  >{\centering\arraybackslash}p{(\linewidth - 4\tabcolsep) * \real{0.2632}}@{}}
\toprule\noalign{}
\begin{minipage}[b]{\linewidth}\raggedright
非零の軸
\end{minipage} & \begin{minipage}[b]{\linewidth}\raggedright
集合形式
\end{minipage} & \begin{minipage}[b]{\linewidth}\centering
状態数
\end{minipage} \\
\midrule\noalign{}
\endhead
\bottomrule\noalign{}
\endlastfoot
\(t\) & \((0, 0, 0, \pm, 0, 0)\) & 2（符号 \(\pm\) の 2 通り） \\
\(R\) & \((0, 0, 0, 0, \pm, 0)\) & 2（符号 \(\pm\) の 2 通り） \\
\(Q\)-遷移 & \((0, 0, 0, 0, 0, Q_\text{初期} \to Q_\text{終期})\) &
9（§3 による） \\
\end{longtable}
}

\textbf{合計}: \(t\) 固定 2 状態 + \(R\) 固定 2 状態 + \(Q\)-遷移 9 状態
= \textbf{13 状態}。

\(xyzt\) のうち \(t\) のみ非零（\(k = 1\)）の場合もボソン型であり、\(n\)
の数え上げにおいて \(t\) が寄与するスピン 1 状態を含む。

\textbf{注記}. \(Q\)-遷移 9 状態からの独立状態の抽出（\(9 \to 8\)
の縮約）、および物理的ゲージボソンとの対応は §9.3--§9.4 で扱う。

\begin{center}\rule{0.5\linewidth}{0.5pt}\end{center}

\subsection{5. スピン 2
集合の数え上げ}\label{ux30b9ux30d4ux30f3-2-ux96c6ux5408ux306eux6570ux3048ux4e0aux3052}

スピン 2 の集合は定義 2.1 によりボソン型かつ \(n = 2\)、すなわち
\(\{t, R, Q\text{-遷移}\}\) のうち 2 本が非零であり、\(xyzt\) の非零軸数
\(k \neq 2\) の状態である。

\(k = 0\)（\(xyzt\) 全軸 \(0\)）の場合、非零となる 2
要素の組合せは以下の 3 通り。

{\def\LTcaptype{none} % do not increment counter
\begin{longtable}[]{@{}ll@{}}
\toprule\noalign{}
非零の 2 要素 & 集合形式 \\
\midrule\noalign{}
\endhead
\bottomrule\noalign{}
\endlastfoot
\(t, R\) & \((0, 0, 0, \pm, \pm, 0)\) \\
\(t, Q\)-遷移 & \((0, 0, 0, \pm, 0, Q_i \to Q_j)\) \\
\(R, Q\)-遷移 & \((0, 0, 0, 0, \pm, Q_i \to Q_j)\) \\
\end{longtable}
}

したがって、本モデルにおけるスピン 2 集合の型は \textbf{3 通り} である。

\begin{center}\rule{0.5\linewidth}{0.5pt}\end{center}

\subsection{6. スピン 3
集合の数え上げ}\label{ux30b9ux30d4ux30f3-3-ux96c6ux5408ux306eux6570ux3048ux4e0aux3052}

スピン 3 の集合は定義 2.1 によりボソン型かつ \(n = 3\)、すなわち
\(\{t, R, Q\text{-遷移}\}\) のすべてが非零であり、\(xyzt\) の非零軸数
\(k \neq 2\) の状態である。

{\def\LTcaptype{none} % do not increment counter
\begin{longtable}[]{@{}ll@{}}
\toprule\noalign{}
非零の 3 要素 & 集合形式 \\
\midrule\noalign{}
\endhead
\bottomrule\noalign{}
\endlastfoot
\(t, R, Q\)-遷移 & \((0, 0, 0, \pm, \pm, Q_i \to Q_j)\) \\
\end{longtable}
}

したがって、本モデルにおけるスピン 3 集合の型は \textbf{1 通り} である。

\begin{center}\rule{0.5\linewidth}{0.5pt}\end{center}

\subsection{7.
非零軸の符号積}\label{ux975eux96f6ux8ef8ux306eux7b26ux53f7ux7a4d}

\textbf{定義 7.1（符号積）}. ボソン集合において、\(t, R\)
軸のうち集合符号が \(0\) でない軸がそれぞれ \(+\) または \(-\)
の符号をとるとき、これらの符号の積を\textbf{符号積} \(P\) と定義する。

\[P = (-1)^{|\{i \in \{t, R\} : \sigma_i < 0\}|}\]

すなわち、\(t, R\) 軸のうち負の符号をとる軸の本数が偶数なら
\(P = +1\)、奇数なら \(P = -1\) である。

\textbf{命題 7.1（各スピン型の \(P\) の分類）}. 各ボソン型の \(P\)
のとりうる値は以下の通り。

{\def\LTcaptype{none} % do not increment counter
\begin{longtable}[]{@{}lll@{}}
\toprule\noalign{}
スピン型 & 非零の \(\{t, R\}\) & \(P\) のとりうる値 \\
\midrule\noalign{}
\endhead
\bottomrule\noalign{}
\endlastfoot
\(s = 0\) & なし & \(P = +1\)（一意） \\
\(s = 1\), \(t\) 型 & \(t\) のみ & \(P = +1\) or \(-1\) \\
\(s = 1\), \(R\) 型 & \(R\) のみ & \(P = +1\) or \(-1\) \\
\(s = 1\), \(Q\)-遷移型 & なし & \(P = +1\)（一意） \\
\(s = 2\), \(tR\) 型 & \(t, R\) 両方 & \(P = +1\) or \(-1\) \\
\(s = 2\), \(tQ\) 型 & \(t\) のみ & \(P = +1\) or \(-1\) \\
\(s = 2\), \(RQ\) 型 & \(R\) のみ & \(P = +1\) or \(-1\) \\
\(s = 3\), \(tRQ\) 型 & \(t, R\) 両方 & \(P = +1\) or \(-1\) \\
\end{longtable}
}

\textbf{証明}. \(t, R\) がともに零のとき \(P = +1\)（負の軸が 0 本）。1
本のみ非零のとき \(\sigma_i = \pm\) により \(P = \pm 1\)。2 本非零のとき
\((+)(+) \to P = +1\), \((+)(-) \to P = -1\), \((-)(+) \to P = -1\),
\((-)(-) \to P = +1\) の 4 配置が可能で \(P = \pm 1\)。\(\square\)

\(P\) の物理的解釈（力の方向性との対応）は §9（解釈の一例）で扱う。

\textbf{注記}. \(R\) 軸を含む正方形面に対する符号付き面積 \(M(\sigma)\)
の配向構造と、それに基づく保存量の厳密な定義は、本稿の範囲を超えるため別論文で扱う。\(M(\sigma)\)
と質量計量の関係についての解釈上の推測は §9.11.1 で述べる。

\begin{center}\rule{0.5\linewidth}{0.5pt}\end{center}

\subsection{8. 本体の結論}\label{ux672cux4f53ux306eux7d50ux8ad6}

本稿では、5 次元超直方体に 4 値離散ラベル軸 \(Q\) を追加した 6
次元構造を定義し、その集合構造から以下の組合せ論的性質を導出した。

\begin{enumerate}
\def\labelenumi{\arabic{enumi}.}
\tightlist
\item
  \textbf{超直方体の要素数}（§1.2）: \(k\) 次元超直方体の \(j\)
  次元面の個数 \(f_j(k) = \binom{k}{j} \cdot 2^{k-j}\)。\(k = 5\)
  の場合、32 頂点、80 辺、80 正方形面、40 立方体胞、10 超直方体。
\item
  \textbf{5 軸の完全対称性}（命題 1.1）: 5
  本の幾何学的軸は組合せ構造上すべて対等であり、任意の軸置換は超直方体の自己同型写像を与える。
\item
  \textbf{軸の表記ラベル}（定義 1.2）: 5 軸に \(x, y, z, t, R\)
  の表記ラベルを与える。
\item
  \textbf{集合符号の厳密な定義}（§1.3.3）: 各軸に対する範囲制約を
  \(+, -, 0, \pm, *\) の 5 記号で表現する。\(Q\) 軸は 4 値（2
  ビット構造）。\(Q = 0\) は \(Q\) 軸に直交する状態を表す。
\item
  \textbf{スピンの定義}（定義 2.1）: \(xyzt\) 4 軸の非零本数 \(k\)
  でフェルミオン型（\(k = 2\)）とボソン型（\(k \neq 2\)）を判定し、\(\{R, Q\text{-遷移}\}\)（フェルミオン型）または
  \(\{t, R, Q\text{-遷移}\}\)（ボソン型）の非零本数 \(n\) からスピンを
  \(s = n + 1/2\) または \(s = n\) と定める。
\item
  \textbf{\(Q\) ラベル間遷移の数え上げ}（§3）: 非零 \(Q\) 値
  \(\{1, 2, 3\}\) 間の遷移は \(3 \times 3 = 9\) 通り。\(Q = 0\)（\(Q\)
  軸に直交）は遷移に関与しない。
\item
  \textbf{スピン 1 集合の数え上げ}（§4）: \(t\) 固定 2 状態 + \(R\) 固定
  2 状態 + \(Q\)-遷移 9 状態 = 13 状態。
\item
  \textbf{スピン 2 集合の 3 型}（§5）: \(tR, tQ, RQ\) の 3 通り。
\item
  \textbf{スピン 3 集合の 1 型}（§6）: \(tRQ\) の 1 通り。
\item
  \textbf{符号積}（定義 7.1、命題 7.1）: \(t, R\) 軸の負の符号の軸数から
  \(P = (-1)^{|\{i : \sigma_i < 0\}|}\) を定義し、各スピン型ごとに \(P\)
  のとりうる値（一意 or 両値）を分類する。
\end{enumerate}

これらはすべて本モデルの外部仮定（5 次元超直方体の採用、\(Q\)
軸の追加）の下で、組合せ論と幾何学から純粋に導出される結果である。

\begin{center}\rule{0.5\linewidth}{0.5pt}\end{center}

\subsection{9.
標準模型との対応の解釈例}\label{ux6a19ux6e96ux6a21ux578bux3068ux306eux5bfeux5fdcux306eux89e3ux91c8ux4f8b}

\textbf{本節の内容は本論文の主張ではなく、解釈の一例である。本モデルから標準模型を物理的に導出することを試みるものではない。}
本節では、§1--§8 で定義・導出した 6
次元構造とその数学的性質が、標準模型の粒子分類とどのように対応付けうるかの一例を示す。

以下の対応は本モデル内部からは正当化されず、\textbf{外部からの解釈の選択}である。将来、別の対応付けが可能になった場合、本節と並列に追加可能である。

\subsubsection{9.1
物理的ラベルの割り当て}\label{ux7269ux7406ux7684ux30e9ux30d9ux30ebux306eux5272ux308aux5f53ux3066}

§1.3.2 で述べた通り、5
軸は組合せ構造上すべて対等であるが、標準模型との対応付けにおいては以下のラベルを割り当てる。

\begin{itemize}
\tightlist
\item
  \(x_1, x_2, x_3 \to xyz\): 3 次元空間軸
\item
  \(x_4 \to t\): 時間軸
\item
  \(x_5 \to R\): 空間スケール軸（主観座標系と 6
  次元構造の空間的スケールを結びつける軸）
\item
  \(Q\): 色荷軸（\(\mathrm{SU}(3)_C\) の離散的符号化）。\(Q = 0\) は
  \(Q\) 軸に直交する状態（色荷なし）を表す。
\end{itemize}

\textbf{\(xyzt\) の初期対称性}: 命題 1.1 により \(xyzt\) の 4
軸は初期状態において完全に対等である。観測においてこのうち 1
軸が「時間」として分離する過程は、自発的対称性の破れとして解釈される。この対称性の破れにより、\(t\)
軸の符号が弱アイソスピンの上下を符号化する（\(t > 0\): 上型、\(t < 0\):
下型、\(t = 0\): 中性）。

\textbf{\(xyz\) の初期対称性と世代}: \(xyz\) の 3
空間軸も、\(\mathrm{SO}(3)\)
対称性のもとで初期状態において完全に対等である。3
世代のフェルミオンは、この初期対称性のもとでは区別されず、すべて同じ質量を持つ。ヒッグスボソンが
\(xyz\) のうち 1 軸に非零成分を取ることで \(\mathrm{SO}(3)\)
対称性が自発的に破れ、結果として世代間の質量差が生じる。本稿では、ヒッグスが非零成分を取る軸を事後的に
\(z\) と表記し、\(z\) 軸と結合するフェルミオンを「第 3
世代」と呼ぶ。世代ラベル（第 1、第 2、第 3）は \(\mathrm{SO}(3)\)
対称性の破れの結果として事後的に定まる概念であり、初期状態では意味を持たない（§9.11.7
参照）。

\subsubsection{\texorpdfstring{9.2 \(Q\) 軸の 4
値と色荷の対応}{9.2 Q 軸の 4 値と色荷の対応}}\label{q-ux8ef8ux306e-4-ux5024ux3068ux8272ux8377ux306eux5bfeux5fdc}

\(Q\) 軸の 4 値は 2 ビット \((c_2, c_3)\)
で構成され、\(\mathrm{SU}(3)_C\)
色荷を符号化する。弱アイソスピンの上下は \(t\) 軸の符号が担い、\(Q\)
軸と \(t\) 軸は完全に独立（直交）な自由度である。

{\def\LTcaptype{none} % do not increment counter
\begin{longtable}[]{@{}ccll@{}}
\toprule\noalign{}
\(Q\) & \((c_2, c_3)\) & 色荷 & 対応 \\
\midrule\noalign{}
\endhead
\bottomrule\noalign{}
\endlastfoot
0 & (0, 0) & 無色（\(Q\) 軸に直交） & レプトン・ボソン \\
1 & (0, 1) & 色1 & クォーク（色1） \\
2 & (1, 0) & 色2 & クォーク（色2） \\
3 & (1, 1) & 色3 & クォーク（色3） \\
\end{longtable}
}

\(Q = 0\) の粒子は \(Q\)
軸に直交する状態であり、色荷を持たない（レプトンおよび全ボソン）。\(Q \neq 0\)
の粒子はクォーク（色荷あり）に対応する。弱アイソスピンの区別は \(t\)
軸の符号が担う（\(t > 0\): 上型、\(t < 0\): 下型、\(t = 0\): 中性）。

反粒子は空間軸（\(x, y, z\)）の符号反転に対応する。\(Q\) 値および \(t\)
軸の符号は変化しない。

\subsubsection{9.3 グルーオン 8
状態への対応}\label{ux30b0ux30ebux30fcux30aaux30f3-8-ux72b6ux614bux3078ux306eux5bfeux5fdc}

§3 で数え上げた \(Q\) 遷移 9 通りのうち、\(Q = 0\)（\(Q\)
軸に直交する色中性モード）に対応する対角成分の 1
つの線形結合（\(\mathrm{SU}(3)\) における色シングレット
\(r\bar{r} + g\bar{g} + b\bar{b}\) に比例する状態）を除いた \textbf{8
独立状態をグルーオン \(g_1, \ldots, g_8\)} に対応させる。

この \(9 - 1 = 8\) の縮約は、\(Q = 0\)
の直交性から自然に理解される。\(Q = 0\) は \(Q\)
軸に直交する状態であり、色荷を運ばないため独立なグルーオンとして数えない。この構造は
\(\mathrm{SU}(3)\) のリー代数 \(\mathfrak{su}(3)\)
の随伴表現（\(3 \otimes \bar{3} = 1 \oplus 8\)
分解）と整合する。\(\mathrm{SU}(3)\) の非可換構造定数 \(f^{abc}\)
は本モデルからは導出されない。

\subsubsection{9.4
ゲージボソンへの対応}\label{ux30b2ux30fcux30b8ux30dcux30bdux30f3ux3078ux306eux5bfeux5fdc}

§4 で数え上げたスピン 1
集合のうち、標準模型のゲージボソンへの対応を以下に示す。

{\def\LTcaptype{none} % do not increment counter
\begin{longtable}[]{@{}llc@{}}
\toprule\noalign{}
固定軸 & 対応 & 状態数 \\
\midrule\noalign{}
\endhead
\bottomrule\noalign{}
\endlastfoot
\(t\) & \(W^+, W^-\) & 2 \\
\(R\) & \(\gamma, Z^0\) & 2 \\
\(Q\)-遷移（§9.3） & \(g_1, \ldots, g_8\) & 8 \\
\textbf{合計} & & \textbf{12} \\
\end{longtable}
}

\(W\) ボソンの 2
状態（\(t = +1, -1\)）は弱アイソスピン遷移（上型⇔下型）を媒介し、\(\gamma\)
と \(Z^0\) は \(R\) 軸の符号（\(+\) vs \(-\)）に対応する。合計 12
は、標準模型のゲージボソン状態数
\(\gamma + g \times 8 + W^+ + W^- + Z^0 = 12\) と一致する。\(W\)
ボソンが \(t\) 軸のみに作用し \(Q\) 軸には作用しないことは、\(Q\) 軸と
\(t\) 軸の直交性と整合する。

\subsubsection{9.5 スピン 2
集合の解釈}\label{ux30b9ux30d4ux30f3-2-ux96c6ux5408ux306eux89e3ux91c8}

§5 で数え上げたスピン 2 集合の 3 型に対して以下の対応を与える。

{\def\LTcaptype{none} % do not increment counter
\begin{longtable}[]{@{}
  >{\raggedright\arraybackslash}p{(\linewidth - 4\tabcolsep) * \real{0.1625}}
  >{\raggedright\arraybackslash}p{(\linewidth - 4\tabcolsep) * \real{0.1750}}
  >{\raggedright\arraybackslash}p{(\linewidth - 4\tabcolsep) * \real{0.6625}}@{}}
\toprule\noalign{}
\begin{minipage}[b]{\linewidth}\raggedright
型
\end{minipage} & \begin{minipage}[b]{\linewidth}\raggedright
固定要素
\end{minipage} & \begin{minipage}[b]{\linewidth}\raggedright
対応
\end{minipage} \\
\midrule\noalign{}
\endhead
\bottomrule\noalign{}
\endlastfoot
\(tR\) 型 & \(t, R\) & \textbf{グラビトン}（標準模型 +
重力の枠組みにおける対応） \\
\(t \cdot Q\)-遷移 型 & \(t, Q\)-遷移 &
標準模型に対応する既知粒子なし \\
\(R \cdot Q\)-遷移 型 & \(R, Q\)-遷移 &
標準模型に対応する既知粒子なし \\
\end{longtable}
}

\(tR\) 型については重力の性質（引力のみ）と整合する。\(tQ\) 型と \(RQ\)
型は、本モデルの組合せ構造から存在が導出される状態であるが、現在の標準模型には対応物がない。

\subsubsection{9.6 スピン 3
集合の解釈}\label{ux30b9ux30d4ux30f3-3-ux96c6ux5408ux306eux89e3ux91c8}

§6 で数え上げたスピン 3 集合の 1
型（\(tRQ\)-遷移）は、標準模型に対応する既知粒子が存在しない。

\subsubsection{9.7
符号積と力の方向性の対応}\label{ux7b26ux53f7ux7a4dux3068ux529bux306eux65b9ux5411ux6027ux306eux5bfeux5fdc}

§7 で定義した符号積 \(P\) に対して、以下の物理的解釈を与える。

\begin{itemize}
\tightlist
\item
  \(P = +1\): 力は一方向的（\textbf{引力のみ}）
\item
  \(P = -1\): 力は双方向的（\textbf{引力・斥力の両方}）
\end{itemize}

命題 7.1 により、スピン 2 の \(tR\) 型は \(P = +1\) と \(P = -1\)
の両配置を許容する。\(P = +1\) の配置（\(t, R\)
が同符号）を重力の媒介と解釈すれば引力のみとなり、重力の性質と整合する。\(P = -1\)
の配置（\(t, R\) が異符号）が物理的に何に対応するかは未解決であり、(a)
物理的に実現されない配置、(b)
未知の現象への対応、のいずれかとして今後の課題である。

同様に、スピン 1 の \(t\) 型・\(R\) 型は \(P = \pm 1\)
の両配置を許容し、引力・斥力の両方を持つ力（電磁力、弱い力）との整合が得られる。スピン
1 の \(Q\)-遷移型は
\(P = +1\)（一意）であり、グルーオンの閉じ込め性との関係が示唆される。

\subsubsection{\texorpdfstring{9.8 \(R\)
軸の値域と力の結合強度}{9.8 R 軸の値域と力の結合強度}}\label{r-ux8ef8ux306eux5024ux57dfux3068ux529bux306eux7d50ux5408ux5f37ux5ea6}

\textbf{\(R\) 軸の値域}: 先行研究 {[}W5{]}
において、離散空間での中心投影の整数論的条件から
\(R = \pm n^2 \cdot L_0\)（\(n \in \mathbb{N}\)、\(L_0\)
は単位長）の離散化が導かれている。\(R\) 軸は主観座標系と 6
次元構造の空間的スケールを結びつける軸であり、質量スペクトルの非有界性に対応する。

\textbf{各力と本モデルの軸の対応の解釈}:

{\def\LTcaptype{none} % do not increment counter
\begin{longtable}[]{@{}
  >{\raggedright\arraybackslash}p{(\linewidth - 6\tabcolsep) * \real{0.0706}}
  >{\raggedright\arraybackslash}p{(\linewidth - 6\tabcolsep) * \real{0.1412}}
  >{\raggedright\arraybackslash}p{(\linewidth - 6\tabcolsep) * \real{0.5529}}
  >{\raggedright\arraybackslash}p{(\linewidth - 6\tabcolsep) * \real{0.2353}}@{}}
\toprule\noalign{}
\begin{minipage}[b]{\linewidth}\raggedright
力
\end{minipage} & \begin{minipage}[b]{\linewidth}\raggedright
媒介粒子
\end{minipage} & \begin{minipage}[b]{\linewidth}\raggedright
関与する軸
\end{minipage} & \begin{minipage}[b]{\linewidth}\raggedright
結合の特徴
\end{minipage} \\
\midrule\noalign{}
\endhead
\bottomrule\noalign{}
\endlastfoot
強い力 & \(g_1\)--\(g_8\) & \(Q\) 軸（色遷移 8 チャンネル, §9.3） &
色荷を持つ粒子間のみ（\(Q \neq 0\)） \\
弱い力 & \(W^\pm, Z^0\) & \(t\) 軸（弱アイソスピン遷移） &
全フェルミオン \\
電磁力 & \(\gamma\) & \(R\) 軸 & 非零電荷を持つ粒子間 \\
重力 & \(G\) & \(t, R\) 2 軸非零（\(n = 2\)） & 全粒子（\(Q\)
非依存） \\
\end{longtable}
}

\subsubsection{9.9 表 A: 確認済み粒子と 6
次元集合の対応}\label{ux8868-a-ux78baux8a8dux6e08ux307fux7c92ux5b50ux3068-6-ux6b21ux5143ux96c6ux5408ux306eux5bfeux5fdc}

本表では粒子のみを示す。反フェルミオンは空間軸（\(x, y, z\)）の符号反転で得られ（\(Q\)
値および \(t\) 軸の符号は不変）、反レプトン 6 + 反クォーク 18 = 24
状態が対応する。ボソンについては \(W^+/W^-\)
のように表中で粒子・反粒子の対を明示しており、\(\gamma, Z^0, H^0, G\),
グルーオンは自己共役である。

\textbf{ボソン}:

{\def\LTcaptype{none} % do not increment counter
\begin{longtable}[]{@{}
  >{\raggedright\arraybackslash}p{(\linewidth - 10\tabcolsep) * \real{0.1408}}
  >{\raggedright\arraybackslash}p{(\linewidth - 10\tabcolsep) * \real{0.1549}}
  >{\centering\arraybackslash}p{(\linewidth - 10\tabcolsep) * \real{0.0704}}
  >{\raggedright\arraybackslash}p{(\linewidth - 10\tabcolsep) * \real{0.4225}}
  >{\centering\arraybackslash}p{(\linewidth - 10\tabcolsep) * \real{0.0704}}
  >{\centering\arraybackslash}p{(\linewidth - 10\tabcolsep) * \real{0.1408}}@{}}
\toprule\noalign{}
\begin{minipage}[b]{\linewidth}\raggedright
粒子
\end{minipage} & \begin{minipage}[b]{\linewidth}\raggedright
記号
\end{minipage} & \begin{minipage}[b]{\linewidth}\centering
\(s\)
\end{minipage} & \begin{minipage}[b]{\linewidth}\raggedright
集合 \((x, y, z, t, R, Q)\)
\end{minipage} & \begin{minipage}[b]{\linewidth}\centering
\(n\)
\end{minipage} & \begin{minipage}[b]{\linewidth}\centering
質量 (MeV)
\end{minipage} \\
\midrule\noalign{}
\endhead
\bottomrule\noalign{}
\endlastfoot
ヒッグス & \(H^0\) & 0 & \((0, 0, +, 0, 0, 0)\) & 0 & 125,250 \\
\(W\) ボソン & \(W^+\) & 1 & \((0, 0, 0, +, 0, 0)\) & 1 & 80,379 \\
\(W\) ボソン & \(W^-\) & 1 & \((0, 0, 0, -, 0, 0)\) & 1 & 80,379 \\
光子 & \(\gamma\) & 1 & \((0, 0, 0, 0, +, 0)\) & 1 & 0 \\
\(Z\) ボソン & \(Z^0\) & 1 & \((0, 0, 0, 0, -, 0)\) & 1 & 91,188 \\
グルーオン & \(g_1\)--\(g_8\) & 1 & \((0, 0, 0, 0, 0, Q_i \to Q_j)\) & 1
& 0 \\
グラビトン & \(G\) & 2 & \((0, 0, 0, \pm, \pm, 0)\) & 2 & 0 \\
\end{longtable}
}

\textbf{検証}: ヒッグス: \(z\) 軸のみ非零で
\(k = 1\)（ボソン型）、\(n = 0\)、\(s = 0\) ✓。\(W\):
\(k = 1\)（ボソン型）、\(t\) が非零で \(n = 1\)、\(s = 1\) ✓。光子:
\(k = 0\)（ボソン型）、\(R\) が非零で \(n = 1\)、\(s = 1\)
✓。グルーオン: \(k = 0\)（ボソン型）、\(Q\)-遷移で \(n = 1\)、\(s = 1\)
✓。グラビトン: \(k = 0\)（ボソン型）、\(t, R\) が非零で
\(n = 2\)、\(s = 2\) ✓。

\textbf{ヒッグスの軸選択と世代}: ヒッグスボソンの非零軸は \(xyz\) の 3
空間軸から 1 本を選ぶ。初期状態では \(xyz\) は \(\mathrm{SO}(3)\)
対称性のもとで対等であるため、この選択は対称性の自発的破れに対応する（§9.1
参照）。ヒッグスが選んだ軸を事後的に \(z\)
と表記し、結合するフェルミオンを「第 3 世代」と呼ぶ。世代番号は
\(\mathrm{SO}(3)\)
対称性の破れの結果として事後的に定まるラベルであり、初期状態では意味を持たない。実測では第
3 世代フェルミオン（トップ: \(172{,}760\)
MeV）が最大質量を持ち、ヒッグスとの結合強度が共有空間軸での重なりに依存するという解釈と整合する（§9.11.7
参照）。

\textbf{レプトン}:

{\def\LTcaptype{none} % do not increment counter
\begin{longtable}[]{@{}
  >{\raggedright\arraybackslash}p{(\linewidth - 12\tabcolsep) * \real{0.2400}}
  >{\raggedright\arraybackslash}p{(\linewidth - 12\tabcolsep) * \real{0.1333}}
  >{\centering\arraybackslash}p{(\linewidth - 12\tabcolsep) * \real{0.0667}}
  >{\centering\arraybackslash}p{(\linewidth - 12\tabcolsep) * \real{0.0667}}
  >{\raggedright\arraybackslash}p{(\linewidth - 12\tabcolsep) * \real{0.2800}}
  >{\centering\arraybackslash}p{(\linewidth - 12\tabcolsep) * \real{0.0667}}
  >{\raggedright\arraybackslash}p{(\linewidth - 12\tabcolsep) * \real{0.1467}}@{}}
\toprule\noalign{}
\begin{minipage}[b]{\linewidth}\raggedright
粒子
\end{minipage} & \begin{minipage}[b]{\linewidth}\raggedright
記号
\end{minipage} & \begin{minipage}[b]{\linewidth}\centering
\(s\)
\end{minipage} & \begin{minipage}[b]{\linewidth}\centering
世代
\end{minipage} & \begin{minipage}[b]{\linewidth}\raggedright
集合
\end{minipage} & \begin{minipage}[b]{\linewidth}\centering
\(Q\)
\end{minipage} & \begin{minipage}[b]{\linewidth}\raggedright
質量 (MeV)
\end{minipage} \\
\midrule\noalign{}
\endhead
\bottomrule\noalign{}
\endlastfoot
電子 & \(e^-\) & 1/2 & 1 & \((+, 0, 0, -, 0, 0)\) & 0 & 0.511 \\
電子ニュートリノ & \(\nu_e\) & 1/2 & 1 & \((+, +, 0, 0, 0, 0)\) & 0 &
\(< 10^{-6}\) \\
ミューオン & \(\mu^-\) & 1/2 & 2 & \((0, +, 0, -, 0, 0)\) & 0 &
105.66 \\
ミューニュートリノ & \(\nu_\mu\) & 1/2 & 2 & \((0, +, +, 0, 0, 0)\) & 0
& \(< 0.17\) \\
タウ & \(\tau^-\) & 1/2 & 3 & \((0, 0, +, -, 0, 0)\) & 0 & 1,776.9 \\
タウニュートリノ & \(\nu_\tau\) & 1/2 & 3 & \((+, 0, +, 0, 0, 0)\) & 0 &
\(< 15.5\) \\
\end{longtable}
}

\textbf{検証}: 電子: \(xyzt\) のうち \(x, t\) が非零で
\(k = 2\)（フェルミオン型）、\(\{R, Q\text{-遷移}\}\)
はともに非該当（\(R = 0\)、\(Q = 0\) は \(Q\)
軸に直交する静的状態で遷移ではない）で \(n = 0\)、\(s = 0 + 1/2 = 1/2\)
✓。ニュートリノ: \(x, y\) が非零で
\(k = 2\)（フェルミオン型）、\(n = 0\)、\(s = 1/2\) ✓。

\textbf{荷電レプトンの特徴}: \(t\)
軸が非零（\(t < 0\)、弱アイソスピン下型）であり、\(Q = 0\)（\(Q\)
軸に直交、無色）。ニュートリノは \(t = 0\)（中性）、\(Q = 0\)
であり、\(xyzt\) のうち空間軸 2 本のみが非零。

\textbf{世代の区別}: 荷電レプトンでは非零の空間軸 1 本と \(t\) 軸で
\(k = 2\) を満たし、ニュートリノでは空間軸 2 本で \(k = 2\)
を満たす。世代ラベルは非零の空間軸によって決定される（世代 1:
\(x\)、世代 2: \(y\)、世代 3: \(z\)）。

\textbf{世代対応規則}: 同世代のレプトン doublet \(\{\ell^-_n, \nu_n\}\)
では、\(\nu_n\) の非零 2 空間軸のうち 1 本が \(\ell^-_n\)
の非零空間軸と一致する。残り 1
軸の選択は解釈上の自由度であり、本稿では巡回的割り当て（世代 1:
\(xy\)、世代 2: \(yz\)、世代 3: \(xz\)）を採用する。

反レプトンは空間軸（\(x, y, z\)）の符号反転に対応する。\(Q\) 値および
\(t\) 軸の符号は変化しない（\(e^+\): \((-, 0, 0, -, 0, 0)\)）。

\textbf{クォーク}:

{\def\LTcaptype{none} % do not increment counter
\begin{longtable}[]{@{}
  >{\raggedright\arraybackslash}p{(\linewidth - 12\tabcolsep) * \real{0.1667}}
  >{\raggedright\arraybackslash}p{(\linewidth - 12\tabcolsep) * \real{0.0667}}
  >{\centering\arraybackslash}p{(\linewidth - 12\tabcolsep) * \real{0.0833}}
  >{\centering\arraybackslash}p{(\linewidth - 12\tabcolsep) * \real{0.0833}}
  >{\raggedright\arraybackslash}p{(\linewidth - 12\tabcolsep) * \real{0.3500}}
  >{\centering\arraybackslash}p{(\linewidth - 12\tabcolsep) * \real{0.0833}}
  >{\raggedright\arraybackslash}p{(\linewidth - 12\tabcolsep) * \real{0.1667}}@{}}
\toprule\noalign{}
\begin{minipage}[b]{\linewidth}\raggedright
粒子
\end{minipage} & \begin{minipage}[b]{\linewidth}\raggedright
記号
\end{minipage} & \begin{minipage}[b]{\linewidth}\centering
\(s\)
\end{minipage} & \begin{minipage}[b]{\linewidth}\centering
世代
\end{minipage} & \begin{minipage}[b]{\linewidth}\raggedright
集合
\end{minipage} & \begin{minipage}[b]{\linewidth}\centering
\(Q\)
\end{minipage} & \begin{minipage}[b]{\linewidth}\raggedright
質量 (MeV)
\end{minipage} \\
\midrule\noalign{}
\endhead
\bottomrule\noalign{}
\endlastfoot
アップ & \(u\) & 1/2 & 1 & \((+, 0, 0, +, 0, 1/2/3)\) & 1,2,3 & 2.16 \\
ダウン & \(d\) & 1/2 & 1 & \((+, 0, 0, -, 0, 1/2/3)\) & 1,2,3 & 4.67 \\
チャーム & \(c\) & 1/2 & 2 & \((0, +, 0, +, 0, 1/2/3)\) & 1,2,3 &
1,270 \\
ストレンジ & \(s\) & 1/2 & 2 & \((0, +, 0, -, 0, 1/2/3)\) & 1,2,3 &
93.4 \\
トップ & \(t\) & 1/2 & 3 & \((0, 0, +, +, 0, 1/2/3)\) & 1,2,3 &
172,760 \\
ボトム & \(b\) & 1/2 & 3 & \((0, 0, +, -, 0, 1/2/3)\) & 1,2,3 & 4,180 \\
\end{longtable}
}

\textbf{検証}: アップクォーク: \(xyzt\) のうち \(x, t\) が非零で
\(k = 2\)（フェルミオン型）、\(Q = 1\) は静的ラベル（遷移ではない）で
\(n = 0\)、\(s = 0 + 1/2 = 1/2\) ✓。

\textbf{クォークの特徴}: 全クォークは \(t\)
軸が非零であり、\(Q \in \{1, 2, 3\}\)（色荷あり）。上型クォーク（\(u, c, t\)）は
\(t > 0\)（弱アイソスピン上型）、下型クォーク（\(d, s, b\)）は
\(t < 0\)（弱アイソスピン下型）。\(t\)
軸の符号が弱アイソスピンの上下および電荷の符号構造を反映する。

各クォークは \(Q\) の色値（\(c_2 c_3 \in \{01, 10, 11\}\)）の 3
色状態を持つ。反クォークは空間軸（\(x, y, z\)）の符号反転に対応する。\(Q\)
値および \(t\) 軸の符号は変化しない。

\textbf{集計}:

{\def\LTcaptype{none} % do not increment counter
\begin{longtable}[]{@{}
  >{\raggedright\arraybackslash}p{(\linewidth - 8\tabcolsep) * \real{0.2674}}
  >{\centering\arraybackslash}p{(\linewidth - 8\tabcolsep) * \real{0.0581}}
  >{\centering\arraybackslash}p{(\linewidth - 8\tabcolsep) * \real{0.1860}}
  >{\centering\arraybackslash}p{(\linewidth - 8\tabcolsep) * \real{0.1628}}
  >{\centering\arraybackslash}p{(\linewidth - 8\tabcolsep) * \real{0.3256}}@{}}
\toprule\noalign{}
\begin{minipage}[b]{\linewidth}\raggedright
分類
\end{minipage} & \begin{minipage}[b]{\linewidth}\centering
\(s\)
\end{minipage} & \begin{minipage}[b]{\linewidth}\centering
\(Q\)
\end{minipage} & \begin{minipage}[b]{\linewidth}\centering
本モデル導出数
\end{minipage} & \begin{minipage}[b]{\linewidth}\centering
SM 対応
\end{minipage} \\
\midrule\noalign{}
\endhead
\bottomrule\noalign{}
\endlastfoot
ヒッグス & 0 & 0 & 1 & 1 \\
ゲージボソン (\(t\) 型) & 1 & 0 & 2 & \(W^+, W^-\) \\
ゲージボソン (\(R\) 型) & 1 & 0 & 2 & \(\gamma, Z^0\) \\
グルーオン (\(Q\)-遷移型) & 1 & 遷移 & 8 & \(g_1, \ldots, g_8\) \\
グラビトン (\(tR\) 型) & 2 & 0 & 1 & \(G\)（SM 外） \\
レプトン & 1/2 & 0 & 6 & 6 \\
反レプトン & 1/2 & 0 & 6 & 6 \\
クォーク & 1/2 & 1,2,3 & 18 & 18 \\
反クォーク & 1/2 & 1,2,3 & 18 & 18 \\
\textbf{小計} & & & \textbf{62} & \textbf{62（61 + G）} \\
\end{longtable}
}

62 状態は標準模型 61 状態 + グラビトン 1 状態にマッピングされる。\(Q\)
軸の 4 値（2 ビット）構造により、\(Q = 0\)（\(Q\)
軸に直交する無色状態）と \(Q \in \{1, 2, 3\}\)（3
色状態）の分離が自然に実現される。

上記に加え、\(s = 3/2, 2, 5/2, 3\) の組合せ（表
B、§9.10）が本モデルから導出されるが、標準模型に対応する既知粒子は存在しない。

\subsubsection{9.10 表 B:
標準模型に対応物のない組合せ}\label{ux8868-b-ux6a19ux6e96ux6a21ux578bux306bux5bfeux5fdcux7269ux306eux306aux3044ux7d44ux5408ux305b}

本モデルの組合せ構造から導出されるが、標準模型に対応する既知粒子が存在しない状態を以下に列挙する。これらは本モデル内の組合せ論的帰結であり、未発見粒子の存在を主張するものではない。

\textbf{ボソン（\(xyzt\) の非零軸数 \(k \neq 2\)）}:

{\def\LTcaptype{none} % do not increment counter
\begin{longtable}[]{@{}
  >{\raggedright\arraybackslash}p{(\linewidth - 8\tabcolsep) * \real{0.1039}}
  >{\centering\arraybackslash}p{(\linewidth - 8\tabcolsep) * \real{0.0649}}
  >{\centering\arraybackslash}p{(\linewidth - 8\tabcolsep) * \real{0.0649}}
  >{\raggedright\arraybackslash}p{(\linewidth - 8\tabcolsep) * \real{0.5195}}
  >{\centering\arraybackslash}p{(\linewidth - 8\tabcolsep) * \real{0.2468}}@{}}
\toprule\noalign{}
\begin{minipage}[b]{\linewidth}\raggedright
型
\end{minipage} & \begin{minipage}[b]{\linewidth}\centering
\(s\)
\end{minipage} & \begin{minipage}[b]{\linewidth}\centering
\(n\)
\end{minipage} & \begin{minipage}[b]{\linewidth}\raggedright
集合形式
\end{minipage} & \begin{minipage}[b]{\linewidth}\centering
状態数
\end{minipage} \\
\midrule\noalign{}
\endhead
\bottomrule\noalign{}
\endlastfoot
\(tQ\) 型 & 2 & 2 & \((0, 0, 0, \pm, 0, Q_i \to Q_j)\) &
\(2 \times N_Q\) \\
\(RQ\) 型 & 2 & 2 & \((0, 0, 0, 0, \pm, Q_i \to Q_j)\) &
\(2 \times N_Q\) \\
\(tRQ\) 型 & 3 & 3 & \((0, 0, 0, \pm, \pm, Q_i \to Q_j)\) &
\(4 \times N_Q\) \\
\end{longtable}
}

ここで \(N_Q\) は \(Q\) 遷移の状態数であり、非零 \(Q\) 値
\(\{1, 2, 3\}\) 間の遷移 \(3 \times 3 = 9\) 通りである（§3）。§9.3
の解釈で色シングレットを除く縮約を行う場合は \(N_Q = 8\) となる。

\textbf{フェルミオン（\(xyzt\) の非零軸数 \(k = 2\)、ただし表 A
に含まれない組合せ）}:

\(xyzt\) から 2 軸を選ぶ \(\binom{4}{2} = 6\) 通り × 各軸の符号
\(2^2 = 4\) 通り × \(Q\) 値 4 通り = \(6 \times 4 \times 4 = 96\)
個のフェルミオン状態が存在する。このうち表 A に対応付けられた 48
状態（レプトン 6 + 反レプトン 6 + クォーク 18 + 反クォーク
18）を除いた残りが未対応状態である。

高スピンフェルミオン（\(R\) 軸または \(Q\)-遷移が非零で \(n \geq 1\)）:

{\def\LTcaptype{none} % do not increment counter
\begin{longtable}[]{@{}lccl@{}}
\toprule\noalign{}
型 & \(s\) & \(n\) & 非零の追加軸 \\
\midrule\noalign{}
\endhead
\bottomrule\noalign{}
\endlastfoot
\(R\) 型 & 3/2 & 1 & \(R\) \\
\(Q\)-遷移型 & 3/2 & 1 & \(Q\)（遷移） \\
\(RQ\) 型 & 5/2 & 2 & \(R, Q\)（遷移） \\
\end{longtable}
}

これらの高スピンフェルミオンは \(xyzt\) の 2 軸が非零に加えて \(R\) や
\(Q\)-遷移が非零となる状態であり、超対称性理論のパートナー粒子とは異なる起源を持つ。

\subsubsection{9.11
解釈例における新たな構造的帰結と未解決問題}\label{ux89e3ux91c8ux4f8bux306bux304aux3051ux308bux65b0ux305fux306aux69cbux9020ux7684ux5e30ux7d50ux3068ux672aux89e3ux6c7aux554fux984c}

\textbf{9.11.1 \(xyzt\) 軸の初期対称性と値の構造}

\(xyzt\) の 4
軸は初期対称性において完全対等であるから、各軸は同種の値（正の整数値）をとることが期待される。ただし、非零軸の本数に応じて
4 種類の値 \(v_1, v_2, v_3, v_4\)（それぞれ非零軸 1 本、2 本、3 本、4
本の場合）が存在し、これらは一般に異なる。

構造的に確定しているのは、\(v_1 \neq v_2 \neq v_3 \neq v_4\)
が許容されるという点のみである。以下はすべて\textbf{解釈上の推測}であり、本モデル内部からは導出されない。

\textbf{次元解析との関係}: 先行研究 {[}W5{]}
の中心投影モデルにおいて、空間的長さ \(L\) と質量 \(M\) は
\(M = L\)（自然単位系）の次元関係を持つ。この関係を本モデルに適用すると、\(xyzt\)
軸の座標値の絶対値の合計が質量計量 \(M(\sigma)\)
に関係する可能性がある。すなわち、集合 \(\sigma = (x, y, z, t, R, Q)\)
に対して

\[M(\sigma) \sim \sum_{i \in \{x, y, z, t\}} |x_i|\]

が質量の起源となりうる。ただし、\(R\) 軸（空間スケール軸）および \(Q\)
軸（色荷軸）がこの計量にどのように関与するかは未確定であり、

\[M(\sigma) \sim f\!\left(\sum |x_i|,\, R,\, Q\right)\]

のような \(R, Q\) を含むより一般的な形式である可能性も排除されない。

\textbf{\(v_1 \gg v_2\) の推測}: ヒッグスボソン（\(k = 1\)、\(z\)
軸のみ非零、\(v_1\)）の質量が \(125{,}250\) MeV
であるのに対し、フェルミオン（\(k = 2\)、非零軸 2
本、\(v_2\)）の大半は遥かに軽い（電子 \(0.511\) MeV、ニュートリノは
\(< 10^{-6}\) MeV）。上記の次元解析的推測が正しいならば、\(v_1 \gg v_2\)
でなければならず、1 軸のみ非零の場合の座標値は 2
軸非零の場合よりも大幅に大きいことが示唆される。

\textbf{世代間質量階層}: ヒッグスボソンが \(xyz\) のうち 1
軸に非零成分を取ることで \(\mathrm{SO}(3)\)
対称性が自発的に破れ、その軸と結合するフェルミオンが最大質量を獲得する。この「選ばれた軸」を事後的に
\(z\) と表記し、結合するフェルミオンを「第 3 世代」と呼ぶ（§9.1
参照）。残る 2 軸（\(x\), \(y\)）に対応する第 1、第 2
世代は、ヒッグスの主振動軸を共有しないため間接的な結合にとどまり、質量が軽い。これは実測での
\(m_3 \gg m_{1,2}\)（トップ \(172{,}760\) MeV \(\gg\) チャーム
\(1{,}270\) MeV、アップ \(2.16\)
MeV）の階層構造と整合する。ただし、\(x\) 軸と \(y\) 軸は
\(\mathrm{SO}(3)\) 対称性により局所的に等価であるため、\(m_2 \neq m_1\)
の差を本枠組み内で導出することはできない。質量比の定量的導出は未解決である。

\textbf{\(k = 3, 4\) の未発見問題との関係}: \(xyzt\) のうち 3 本または 4
本が非零（\(k = 3, 4\)）の状態は、ボソン型（\(k \neq 2\)）として本モデルから導出されるが、表
A に対応する既知粒子は存在しない。これらの状態に対応する値 \(v_3, v_4\)
が極めて小さい（あるいは零に近い）場合、対応する粒子のエネルギースケールが現在の実験的到達範囲を超える可能性がある。逆に
\(v_3, v_4\)
が極めて大きい場合は、安定性の問題から物理的に実現されない可能性もある。いずれにせよ、\(k = 3, 4\)
の状態が未発見であることと \(xyzt\)
軸の値の構造は密接に関連している可能性がある。

\textbf{9.11.2 \(Q = 0\) 制約と自由粒子}

\(Q = 0\) は \(Q\)
軸に直交する状態であり、色荷を持たないことを意味する。解釈として、\(Q = 0\)
の粒子のみが自由粒子として観測可能であるという制約が考えられる。これはクォークの閉じ込め（\(Q \neq 0\)
の粒子は単独で観測されない）の幾何学的起源として解釈可能である。

\textbf{9.11.3 未発見粒子の予言}

本モデルの組合せ構造から導出される状態のうち、表 A の 62
状態を超えるものは、(a) 未発見粒子の予言、(b)
物理的に実現されない組合せ（選択則の存在）、(c)
本解釈の限界、のいずれかとして今後の研究で検討される。

特に、高スピンフェルミオン（\(s = 3/2, 5/2\)）および追加ボソン（\(tQ\)
型、\(RQ\) 型のスピン 2、\(tRQ\) 型のスピン
3）は、既知粒子に対応物を持たない。

\textbf{9.11.4 質量比の導出}

本モデルの集合構造から既知の粒子質量比を導出することは今後の課題である。§9.11.1
で述べた \(xyzt\) 軸の 4 種の値 \(v_1, v_2, v_3, v_4\)
の具体的な比、および \(R\) 軸・\(Q\) 軸の質量計量 \(M(\sigma)\)
への関与の形式が、この問題の鍵となる可能性がある。特に \(v_1 \gg v_2\)
の関係（§9.11.1）が確認されれば、ヒッグスボソンとフェルミオンの質量階層に対する幾何学的説明が得られる。

\textbf{9.11.5 ゲージ群の導出}

\(Q\) 軸の 4 値（2 ビット構造）と \(\mathrm{SU}(3)_C\)、\(t\) 軸の符号と
\(\mathrm{SU}(2)_L\)、\(R\) 軸の連続スケールと \(\mathrm{U}(1)_Y\)
の間の対応が示唆されるが、離散構造から連続的ゲージ群への移行の厳密な機構は未証明である。

\textbf{9.11.6 CKM 行列と PMNS 行列}

世代間の混合（CKM 行列、PMNS 行列）が、\(xyz\) 軸の \(S_3\)
対称性の破れとしてどのように導出されるかは未解決である。

\textbf{9.11.7 ヒッグス機構と世代質量階層}

ヒッグスボソンは \(xyz\) 3 軸のうち 1 軸に非零成分を取る（\(k = 1\)、表
A: \((0, 0, +, 0, 0, 0)\)）。§9.1 で述べた通り、初期状態では \(xyz\) は
\(\mathrm{SO}(3)\)
対称性のもとで完全に対等であるため、どの軸を選ぶかは対称性の自発的破れによる事後的な選択である。本稿では、ヒッグスが選ぶ軸を
\(z\) と表記する。

ヒッグスの \(z\) 軸選択により:

\begin{itemize}
\tightlist
\item
  \(z\)
  軸を非零に持つフェルミオンはヒッグスとの結合が最大となり、最大質量を獲得する（これを「第
  3 世代」と呼ぶ。実測ではトップ: \(172{,}760\) MeV 等）
\item
  \(x\) 軸、\(y\)
  軸を非零に持つフェルミオンはヒッグスの主振動軸を共有せず、結合が弱いため質量が軽い（第
  1、第 2 世代）
\end{itemize}

このように、世代番号と質量階層 \(m_3 \gg m_{1,2}\)
は、初期対称性の自発的破れの結果として事後的に現れる。ヒッグス場による自発的対称性の破れの力学的機構が本枠組みでどう理解されるかは今後の課題である。

\textbf{9.11.8 電荷値の完全な導出}

\(t\) 軸の符号から電荷の正負が、\(Q\) 軸の \(c_2 c_3\)
から色荷の有無と電荷の絶対値（1 vs 1/3）が読み取れるが、\(+2/3\) と
\(-1/3\) の非対称性の完全な導出は今後の課題である。

\subsubsection{9.12
解釈例の結論}\label{ux89e3ux91c8ux4f8bux306eux7d50ux8ad6}

本節では、§1--§8 で定義・導出した 6
次元構造の数学的性質を、標準模型の粒子分類と対応付ける一例を示した。

\begin{enumerate}
\def\labelenumi{\arabic{enumi}.}
\tightlist
\item
  \(Q\) 軸の 4 値（2 ビット構造）が \(\mathrm{SU}(3)_C\)
  色荷を離散的に符号化し、\(t\) 軸の符号が \(\mathrm{SU}(2)_L\)
  弱アイソスピンを符号化する。\(Q\) 軸と \(t\)
  軸は完全に独立（直交）な自由度である
\item
  \(xyzt\) 4
  軸の初期対称性から、スピンの定義が自然に導かれる（フェルミオン:
  \(k = 2\)、ボソン: \(k \neq 2\)）
\item
  62 状態が標準模型の既知粒子（61 + グラビトン）にマッピング可能
\item
  \(Q\) 遷移 9 通りのうち \(Q = 0\)（\(Q\)
  軸に直交する色シングレット）を除いた 8
  状態がグルーオンに対応（\(9 \to 8\) の縮約は \(\mathrm{SU}(3)\)
  表現論に依存）
\item
  スピン 1 集合の 12 状態が標準模型のゲージボソンに対応
\item
  荷電フェルミオンは \(t\) 軸が非零であり、\(t\)
  の符号が弱アイソスピンの上下（上型: \(t > 0\)、下型:
  \(t < 0\)）を反映する
\item
  \(Q = 0\)（\(Q\)
  軸に直交する状態）制約がクォーク閉じ込めの幾何学的起源として解釈可能
\item
  \(xyzt\) 軸の値の構造（4
  種の正整数値）が質量階層の起源を含む可能性がある
\item
  高スピン状態および組合せ論的余剰状態は未発見粒子の予言または選択則の存在を示唆する
\end{enumerate}

具体的な質量値、電荷値、ゲージ群の非可換構造、世代混合行列の導出は今後の課題である（§9.11）。

\begin{center}\rule{0.5\linewidth}{0.5pt}\end{center}

\subsection{参考文献}\label{ux53c2ux8003ux6587ux732e}

{[}1{]} Finkelstein, D. \& Rubinstein, J. ``Connection between Spin,
Statistics, and Kinks.'' \emph{J. Math. Phys.} 9, 1762 (1968).

{[}2{]} Atiyah, M.F., Manton, N.S. \& Schroers, B.J. ``Geometric Models
of Matter.'' \emph{Proc. R. Soc. A} 468, 1252 (2012). arXiv: 1108.5151.

{[}3{]} Furey, C. ``Standard Model Physics from an Algebra?'' arXiv:
1611.09182 (2016).

{[}4{]} Kaluza, Th. ``Zum Unitätsproblem der Physik.''
\emph{Sitzungsber. Preuss. Akad. Wiss. Berlin} 966--972 (1921).

{[}5{]} Coxeter, H.S.M. \emph{Regular Polytopes}. Dover (1973).

{[}6{]} Kihara, N. ``Bivectorial Classification of Spin Derived from the
Orientation Structure of the 5-Dimensional Hypercube.'' DOI:
\href{https://doi.org/10.5281/zenodo.19643358}{10.5281/zenodo.19643358}
(2026).

{[}7{]} Kihara, N. ``Geometric Classification of Spin Derived from the
Orientation Structure of the 5-Dimensional Orthoplex.'' DOI:
\href{https://doi.org/10.5281/zenodo.19630972}{10.5281/zenodo.19630972}
(2026).

{[}W5{]} Kihara, N. ``Complete Spherical Coverage of Central Projection
in Discrete Space and the Number-Theoretic Necessity of the
5-Dimensional Background Space.'' DOI:
\href{https://doi.org/10.5281/zenodo.19624957}{10.5281/zenodo.19624957}
(2026).
\end{document}
